% --- Section: Introduction ---
% Drafted Phase 5 — values pulled from cleanup/04_phase5_data_pulled.md
% Target: ~600–800 words.

\section{Introduction}
\label{sec:intro}

Many users of large language models pick a model by reading the top
of a leaderboard. For statistical reasoning that habit is risky.
Benchmarks such as StatEval~\cite{lu2025stateval},
MathEval~\cite{liu2025matheval}, the human-comparison study of
Nagarkar et al.~\cite{nagarkar2026canllm}, and ad-hoc evaluations on
Bayesian tasks~\cite{longjohn2025bayesian} report accuracy alone.
But accuracy fuses three properties that a downstream user actually
cares about separately: whether the model lands on the right number,
whether it does so when the prompt is reworded or the numbers are
swapped, and whether its stated confidence tracks its empirical
correctness. When those three disagree, a single accuracy column is
misleading.

We make the disagreement concrete on 171 Bayesian and inferential
tasks. Five frontier LLMs are scored on accuracy, robustness, and
calibration; the rank orders differ. Two patterns are doing most of
the work. GPT-4.1, last on accuracy (0.621), is first on robustness
(0.9996); its base and perturbation pass rates are within 0.0003 of
each other. Claude leads accuracy (0.683) and binary-pass calibration
(ECE 0.0334) but drops to fourth on robustness (0.9565), losing three
percentage points between base and perturbed prompts. A consumer who
reads only the accuracy column never sees either pattern, and the
deployment advice each one implies is different. Separately, 44.3\%
of responses fail to state the statistical assumptions the derivation
requires, even when the final number is correct, and assumption
violation is the dominant failure mode at 47\% of judge-classified
failures. Du et al.~\cite{du2025icecream} report the same blindness
in causal reasoning.

The methodological gap behind the headline also bears naming. Most
LLM-statistics evaluations rely on a keyword rubric, a single judge,
or both, and both have known biases. Keyword rubrics confuse surface
form with substance: the word ``assume'' in a response is not the
same thing as actually stating an assumption. Single-judge protocols
suffer self-preference bias and design-choice sensitivity, as
documented by Yamauchi et al.~\cite{park2025judge} and Feuer et
al.~\cite{judgment2025noise}. We score each response twice, once with
a deterministic keyword rubric in the live runner and once with an
external, out-of-family judge: Llama 3.3 70B Turbo via Together AI,
not one of the five benchmarked models. The two views disagree where
it matters. They agree only moderately on assumption-compliance
(Spearman 0.602) and barely at all on method structure and reasoning
quality. The keyword rubric also passes assumptions more leniently
than the judge (71.1\% vs 55.7\%), which is why keyword-only
protocols read Bayesian competence higher than it is.

Our contributions are:
\begin{enumerate}
  \item A benchmark of 171 Bayesian and inferential tasks (136
        closed-form conjugate models and 35 computational Bayesian
        methods covering Gibbs, Metropolis-Hastings, HMC, RJMCMC,
        hierarchical, VB, and ABC) with 473 perturbations along
        three axes (rephrase, semantic, numerical) that touch all
        38 task types.
  \item An external-judge protocol using an out-of-family open-weights
        model (Llama 3.3 70B Turbo) for assumption-compliance and
        reasoning quality. Keyword-rubric vs judge agreement is
        reported per dimension (Spearman 0.602 on
        assumption-compliance, much lower elsewhere). The split
        between what a rubric can see and what a judge can see is
        itself a finding.
  \item Empirical evidence that accuracy, robustness, and calibration
        rankings of the same five LLMs disagree non-trivially. The
        two flagship divergences are a 5-to-1 inversion between
        accuracy and robustness (GPT-4.1) and a 1, 1, 4 pattern for
        the accuracy-and-calibration leader (Claude). Neither shows
        up on a single-metric leaderboard.
  \item An open audit trail. NMACR scoring weights (A=0.30, R=0.25,
        M=0.20, C=0.15, N=0.10), aggregation choices, and tolerance
        thresholds are all documented and tied to the papers that
        motivated them.
\end{enumerate}

Section~\ref{sec:related} reviews statistical-reasoning benchmarks,
the calibration literature~\cite{guo2017calibration}, and
LLM-as-judge work~\cite{zheng2023judge,park2025judge,judgment2025noise}.
Section~\ref{sec:methodology} covers the task suite, perturbations,
judge protocol, and experimental setup. Section~\ref{sec:results}
presents the three rankings. Section~\ref{sec:discussion} argues
the divergence is not noise. Section~\ref{sec:conclusion} closes
with threats to validity and future work.
